\chapter{Discussion}
\label{ch:discussion}
\glsunset{ADNEX}
\glsunset{RMI}
\glsunset{IOTA}
\glsunset{AUC}
\glsunset{EIDM}
\glsunset{IPD}
\glsunset{SaMD}
\glsunset{MDR}
\glsunset{TRIPOD}
\glsunset{PROBAST}


This thesis set out to improve evidence-informed decision-making in the diagnosis of ovarian cancer by evaluating the \gls{ADNEX} model and advancing statistical methodology for clinical prediction models. The work integrates an applied arm --- focused on evaluating, comparing and updating \gls{ADNEX} in diverse clinical settings --- with a methodological arm that advances methodology in model development, evaluation and probability estimation. 

\section{Summary of findings}

Chapter~\ref{sec:chp_ADNEXSR} presented the first comprehensive systematic review and meta-analysis of \gls{ADNEX} external validation studies. Across 47 studies encompassing more than 17\,000 tumours, the model demonstrated robust discriminative performance, with a summary \gls{AUC} of 0.93 for distinguishing benign from malignant masses \cite{barrenadaADNEXRiskPrediction2024}. However, the review also exposed important weaknesses in the validation literature: 91\% of validations were at high risk of bias, primarily due to small sample sizes, inadequate handling of missing data, and incomplete performance assessment. Perhaps most critically, only 4 of 47 studies reported any form of calibration assessment, making it impossible to draw conclusions about the accuracy of \gls{ADNEX}'s predicted probabilities across settings.

Chapter~\ref{sec:chp_ADNEXvsRMI} extended this work with a head-to-head meta-analysis comparing \gls{ADNEX} to the \gls{RMI}, the most widely used alternative in clinical practice. The results were unequivocal: \gls{ADNEX} consistently outperformed \gls{RMI} in terms of discrimination (summary \gls{AUC} 0.92 vs 0.85 in operated patients) and clinical utility. The probability of \gls{ADNEX} being clinically useful in a hypothetical new centre was 96\%, compared to only 15\% for \gls{RMI} at the 10\% risk threshold \cite{barrenadaHeadHeadComparison2025}. These findings provide strong evidence that clinical guidelines should recommend \gls{ADNEX} over \gls{RMI} for the preoperative characterisation of ovarian masses \cite{sundarRiskpredictionModelsPostmenopausal2024,sundarDiagnosticTestsOvarian2026}.

A recurring theme throughout this thesis is that calibration --- the agreement between predicted probabilities and observed event rates --- remains underappreciated in the clinical prediction literature. The systematic reviews in Chapters~\ref{sec:chp_ADNEXSR} and~\ref{sec:chp_ADNEXvsRMI} showed that most validation studies of \gls{ADNEX} reported discrimination and classification but ignored calibration. That imbalance is methodologically weak and clinically risky. A model can rank patients correctly and still provide systematically wrong absolute probabilities, which can be harmful when the output is used to decide whether a woman with an adnexal mass should be referred, conservatively managed, or operated on \cite{vancalsterCalibrationRiskPrediction2015}. In the context of \gls{ADNEX} and prediction models as decision support tools, calibration and the probability estimate is not a secondary output: it is the output that informs action.

The same applied story motivates Chapter~\ref{sec:chp_ADNEX2}. The original \gls{ADNEX} development data were necessarily limited to women who underwent surgery, whereas in clinical practice the model is intended to be used to determine who to operate. We used bayesian refitting borrowing information from the original model to update \gls{ADNEX} with a larger and more representative dataset that includes non-operated patients. The goal is to produce  the \emph{ADNEX2} model that is better calibrated and more robust for use in a two-step strategy that starts with the simple benign descriptors, followed by the full model for indeterminate cases. ADNEX2 is undergoing medical device regulatory approval, and will be deployed in the near future through Gynaia Inc. (https://www.gynaia.com/) in a safe and easy to use app. 

Chapter~\ref{sec:chp_RFOverfitting} shifted attention from the applied model to the behaviour of algorithms used for probability estimation. Using both a case study in ovarian malignancy prediction and simulation experiments, the chapter showed that random forests can achieve almost perfect training \gls{AUC}s, which is often a sign of overfitting, and at the same time competitive test performance. The key issue is not classical overfitting in the narrow sense, but the tendency of fully grown trees to create local probability peaks and unstable risk estimates. That distinction matters because prediction models are often tuned as if classification performance were the only objective. For risk prediction, however, the quality of the probabilities themselves is central, and the default strategy of growing trees as deep as possible can be counterproductive when the goal is calibration rather than only discrimination. 

Chapter~\ref{sec:chp_ClusteredCalibration} addressed a different but closely related issue: calibration should not be treated as a purely pooled property when the data come from multiple centres. Much of the evidence used to evaluate clinical prediction models is clustered by hospital or study, yet traditional calibration plots suppress that structure. The chapter therefore proposed clustered group calibration (CG-C), two-stage meta-analysis calibration (2MA-C), and mixed model calibration (MIX-C). The main lesson is that calibration has both an overall component and a between-centre component. The recommended approaches, especially 2MA-C with splines and MIX-C, make that distinction visible by combining an overall calibration curve with prediction intervals where the calibration of a new centre should lie. That is a more realistic summary than a single pooled curve because it asks the clinically relevant question of whether a model remains well calibrated when it is moved to a new setting.

Chapter~\ref{sec:chp_FundamentalProblem} took a step back from the technical evaluation of model performance to examine a more fundamental question: how much can we trust a risk estimate for an individual patient? The study demonstrated that individual risk estimates are far more uncertain than commonly assumed. While estimation uncertainty (due to finite sample size) decreases with larger training samples, model uncertainty (due to arbitrary modelling choices) and applicability uncertainty (due to differences in measurement procedures and populations) remain substantial. Even with 10\,000 training observations, the 95\% range of predicted probabilities for the same patient averaged 0.39 on the probability scale. These findings challenge the common practice of presenting a single risk estimate to patients and clinicians without acknowledging the underlying uncertainty.

Chapter~\ref{sec:chp_varselection} is still under development, but it follows naturally from the methodological chapters. If a model is judged not only by fit but also by usefulness, then predictor selection should account for the clinical consequences of including or excluding variables. That question is especially relevant in settings such as ovarian mass diagnosis, where a model may guide referral, surgery, or follow-up decisions.



\section{Embracing uncertainty in individualised decision-making}

Clinical prediction models are developed and evaluated at the population level, but their purpose is to inform decisions for individual patients \cite{steyerbergClinicalPredictionModels2009}. The chapters in this thesis show that moving from population performance to individual decision-making adds successive layers of uncertainty. Chapter~\ref{sec:chp_ADNEXSR} established that \gls{ADNEX} performs well on average across settings. Chapter~\ref{sec:chp_ADNEXvsRMI} showed that it compares favourably with \gls{RMI} for triage. Chapter~\ref{sec:chp_ClusteredCalibration} showed that pooled performance can mask meaningful centre-to-centre variation. Chapter~\ref{sec:chp_FundamentalProblem} then showed that uncertainty persists even after all of those population-level summaries are taken into account.

That sequence matters because the apparent precision of a risk estimate can be misleading. A patient and clinician may see one number, but that number is the product of a model that may behave differently in another centre, another patient subgroup, or another data collection context. Chapter~\ref{sec:chp_ClusteredCalibration} makes this visible at the centre level; Chapter~\ref{sec:chp_FundamentalProblem} makes it visible at the individual level. 
Chapters~\ref{sec:chp_ADNEXSR} and~\ref{sec:chp_ADNEXvsRMI} established that \gls{ADNEX} performs well on average across diverse populations. However, both chapters noted considerable heterogeneity in performance across centres and studies. Chapter~\ref{sec:chp_ClusteredCalibration} provided formal methods to quantify this heterogeneity, showing that calibration can vary substantially between clusters even when the overall curve appears satisfactory. 
This has direct implications for the concept of \gls{EIDM} introduced earlier in the thesis. A predicted probability should be treated as an estimate based on the used model, not as an exact measurement. Confidence intervals will represent how uncertain the model is based on the used data, but do not guarantee that the true underliying risk of the patient lies within the interval. In practice, that means risk estimates should be communicated together with an explanation of what is known, what is uncertain, and why the estimate may change if the patient were seen in another centre or assessed under slightly different assumptions \cite{pateUncertaintyUsingRisk2019, rileyUncertaintyRiskEstimates2025, aldermanTacklingAlgorithmicBias2025, ganapathiTacklingBiasAI2022}.
Chapter~\ref{sec:chp_FundamentalProblem} pushed this line of reasoning further, showing that even within a single centre, the predicted probability for an individual patient is subject to model and applicability uncertainty that cannot be reduced by increasing sample size alone \cite{pateUncertaintyUsingRisk2019, rileyUncertaintyRiskEstimates2025}. A risk estimate of 15\% should not be interpreted as a precise measurement but rather as a point within a range of plausible values, each of which might lead to a different clinical decision.
These findings do not diminish the value of clinical prediction models. \gls{ADNEX} remains a useful tool for ovarian mass characterisation because it shows clinical utility at population level and can support referral and management decisions. The lesson is instead that its outputs should be used with caution. Clinical prediction models should support shared decision-making rather than replace it, and that requires assesments embracing this uncertainty \cite{moonsCriticalAppraisalData2014, steyerbergClinicalPredictionModels2009}.

\section{Towards honest and safe medical devices}

The systematic reviews in Chapters~\ref{sec:chp_ADNEXSR} and~\ref{sec:chp_ADNEXvsRMI} consistently identified poor methodological quality in external validation studies. This is very concerning because this evidence is of utmost importance for safe implementation of clinical prediction models \cite{steyerbergPredictionModelsNeed2016}.  The most common issues --- assessed using the \gls{PROBAST} tool \cite{wolffPROBASTToolAssess2019} --- were inadequate sample sizes, inappropriate exclusion of patients with missing data, and incomplete reporting of model performance. Adherence to the \gls{TRIPOD} reporting guidelines \cite{collinsTRIPODAIStatementUpdated2024} was low, with key items such as sample size justification, missing data handling, and calibration assessment reported in fewer than 8\% of studies. These findings are consistent with previous reviews of clinical prediction model validation studies, which have repeatedly highlighted similar issues \cite{collinsExternalValidationMultivariable2014, bouwmeesterReportingMethodsClinical2012, andaurnavarroCompletenessReportingClinical2022,andaurnavarroSystematicReviewFinds2023, andaurnavarroSystematicReviewIdentifies2023, dhimanMethodologicalConductPrognostic2022, dhimanOverinterpretationFindingsMachine2023}.

These findings are even more concerning in light of the regulatory developments discussed in the introduction. As clinical prediction models increasingly fall under medical device regulation --- including the EU \gls{MDR} and the \gls{SaMD} framework --- the evidence base for their validation must improve. Regulatory frameworks require demonstration of safety and effectiveness, which in the context of prediction models translates to robust evidence of discrimination, calibration, and clinical utility across diverse populations \cite{vancalsterPerformanceEvaluationPredictive2024, collinsTRIPODAIStatementUpdated2024}. The current state of the validation literature and regulatory approval, falls short of what should be expected for tools that are intended to be used in patients. Recent studies have highlighted that many models that have been approved and are in use, have not undergone rigorous external validation, and their performance in real-world settings is often unknown or unclear \cite{mehtaEvaluatingTransparencyAI2025,vanleeuwenArtificialIntelligenceRadiology2021}.

Improving this situation requires combined efforts from researchers and regulators. Researchers should use better study designs, with adequate sample size\cite{rileyMinimumSampleSize2021}, proper handling of missing data and prospective multicenter data collection. The quality of reporting in published literature also needs improving which can be done by adhering to international guidelines like \gls(TRIPOD) and its extensions \cite{collinsTRIPODAIStatementUpdated2024,debrayTransparentReportingMultivariable2023, snellTransparentReportingMultivariable2023}), and a shift in culture towards recognising external validation as an ongoing process rather than a one-time exercise is needed \cite{vancalsterThereNoSuch2023}. From the regulators side, there should be stricter requirements about the model's performance before approval, ensuring that at least key performance measures are presented such as the calibration plot and a clinical utility metric (i.e. Net Benefit) \cite{vancalsterPerformanceEvaluationPredictive2024, ShowUsEvidence2026}. Both perspectives are synergistic: better research leads to better evidence for regulators, and stricter regulation incentivises better research. Ultimately, the goal is to ensure that clinical prediction models are safe, effective, and trustworthy tools for improving patient care.


\section{Future directions for trustworthy AI in medicine}

There is an exponential growth in the development of clinical prediction models, but the findings of this thesis suggest that there is still much work to be done to ensure that these models are trustworthy and useful in practice. A recent systematic review showed that the probability of a new clinical prediction model being externally validated in 5 years is of 13\% \cite{arshiNumberPublicationsNew2025}. Only one in eight models will be validated, which is a requirement for being trustworthy and applied in clinical practice. The findings of this thesis suggest that there is need for more and better validation studies, with a focus on calibration and clinical utility. This is a fast growing and evolving field, and there are many opportunities for future research to improve the development, evaluation, and implementation of clinical prediction models. Some of the most promising directions include:

\subsection{Focus on clinical utility}

\subsection{Better uncertainty estimation}

\subsection{Evaluation when clustering is present}

\subsection{Policymakers and regulators as key stakeholders}

\subsection{Trustworthy AI for ovarian cancer diagnosis}

\subsection{Updating ADNEX}
The most immediate applied priority is the development and validation of \gls{ADNEX}2. The preceding chapters make clear what such an update should address: the development data should be expanded to include non-operated patients, more recent \gls{IOTA} data should be incorporated, and the updated model should be validated with methods that explicitly quantify calibration in clustered settings \cite{vancalsterEvaluatingRiskOvarian2014, steyerbergAssessmentHeterogeneityIndividual2019}. The aim is not simply to improve average performance, but to obtain a model whose probabilities remain reliable across centres and clinical contexts.

\subsection{Variable selection for clinical utility}
Current variable selection methods in prediction modelling typically optimise statistical criteria such as the Akaike Information Criterion, cross-validated log-loss, or discrimination. However, the ultimate purpose of a diagnostic prediction model is to improve clinical decisions. Developing variable selection algorithms that directly optimise for net benefit \cite{vickersNetBenefitApproaches2016} or related utility measures could lead to simpler models that are more useful in practice, even if they sacrifice a small amount of traditional statistical performance. This is the focus of the ongoing work in Chapter~\ref{sec:chp_varselection}.

\subsection{Extending clustered calibration methodology}
The methods proposed in Chapter~\ref{sec:chp_ClusteredCalibration} can be extended in several directions. Bayesian approaches to prediction interval estimation may improve coverage properties compared to the current frequentist methods. More complex heterogeneity structures, including \glspl{random_slope} and nonlinear centre effects, should be evaluated. The methods should also be tested with prediction models beyond logistic regression, including machine learning algorithms whose probability estimates may exhibit different calibration patterns.

\subsection{Uncertainty communication in clinical practice}
Chapter~\ref{sec:chp_FundamentalProblem} demonstrated that individual risk estimates carry substantial uncertainty, but translating this insight into clinical practice remains an open challenge. Future work should explore how uncertainty can be communicated effectively to clinicians and patients without undermining trust in prediction models \cite{kompaSecondOpinionNeeded2021}. Practical approaches may include prediction abstention, presenting ranges instead of point estimates, and integrating uncertainty quantification into clinical decision support systems. These developments connect directly to the broader goal of \gls{EIDM} and shared decision-making.

\subsection{Model lifecycle monitoring}
As clinical prediction models are increasingly recognised as medical devices, continuous monitoring of their performance after deployment becomes a regulatory and ethical requirement. Methods for detecting performance degradation over time --- due to changes in patient populations, clinical practice, or measurement technologies --- need to be developed and standardised. The clustered calibration methods from this thesis provide a foundation for such monitoring, but further work is needed to adapt them for sequential surveillance rather than one-time validation.



\section{Conclusion}
This thesis contributes to the development and evaluation of clinical prediction models for ovarian cancer diagnosis from both applied and methodological perspectives. The systematic evaluation of \gls{ADNEX} confirmed its value as a diagnostic tool and provided the evidence base for recommending it over alternatives such as the \gls{RMI}. At the same time, the methodological chapters showed why this should not be the end of the story: calibration is rarely assessed, probability estimates from machine learning algorithms can be misleading if not validated properly, clustering changes how performance should be interpreted, and individual risk predictions carry more uncertainty than is often acknowledged.

These findings underscore that clinical prediction models, including \gls{ADNEX}, are very valuable but they need to be used responsibly. Trustworthy AI or prediction model requires honest model development and validation, ensuring that the results are reliable and generalisable. Their responsible use requires rigorous validation, transparent communication of uncertainty, and ongoing monitoring --- principles that are increasingly formalised in both reporting guidelines and medical device regulation, although not properly enforced by some regulators \cite{,mehtaEvaluatingTransparencyAI2025}. The ongoing work on variable selection for clinical utility and the \gls{ADNEX}2 update aims to translate the methodological insights of this thesis into a better model for clinical practice.

\color{black}
